\editionPage{indices-defis}{Indications progressives}{Lire une aide, puis reprendre la recherche}{Les indices orientent la démarche ; les solutions restent dans la partie Corrigés.}
\begin{blocEdition}{indices-defis}
\editionEtape{D01 · Mélanger n'est pas moyenner}
\textbf{Indice 1.} Ce ne sont pas les pH qui s'additionnent, mais les quantités d'ions oxonium.
Calculer la quantité de \ce{H3O+} apportée par chaque solution, puis la concentration dans le volume total.\par\smallskip
\textbf{Indice 2.} Pour la question 2, poser $a=10^{-p_1}$ et $b=10^{-p_2}$ et utiliser la décroissance
de la fonction $x\mapsto-\log x$. Pour la question 4, remplacer $[\ce{HO-}]$ grâce à $K_\mathrm{e}$,
puis ne garder que la racine positive du trinôme.\par\smallskip
\editionEtape{D02 · La méthode des ajouts dosés}
\textbf{Indice 1.} Dans chaque fiole, la quantité de X est la somme de celle apportée par l'échantillon et de celle
apportée par l'étalon ; diviser par le volume de la fiole pour obtenir la concentration.\par\smallskip
\textbf{Indice 2.} L'abscisse du point où $A=0$ vaut $-b/a$ ; elle est négative et sa valeur absolue est le volume
d'étalon qui contiendrait autant de X que l'échantillon. Une absorbance parasite ne modifie que l'ordonnée à l'origine.\par\smallskip
\editionEtape{D03 · Identifier un acide inconnu}
\textbf{Indice 1.} Multiplier chaque pourcentage par la masse molaire supposée, puis diviser par la masse molaire
de l'élément : on obtient le nombre d'atomes de chaque sorte dans la molécule.\par\smallskip
\textbf{Indice 2.} Avec $p$ groupes carboxyle, $n(\text{acide})=n(\ce{HO-})/p$. Chaque groupe carboxyle contient une
double liaison C=O : le nombre de doubles liaisons doit être au moins égal à $p$.\par\smallskip
\end{blocEdition}

\editionPage{indices-defis-2}{Indications progressives (suite)}{Lire une aide, puis reprendre la recherche}{Les indices orientent la démarche ; les solutions restent dans la partie Corrigés.}
\begin{blocEdition}{indices-defis-2}
\editionEtape{D04 · Ordre 1 ou ordre 2 ?}
\textbf{Indice 1.} Pour une loi d'ordre 1, la concentration est divisée par deux à chaque durée $t_{1/2}$ : à la date
$2\,t_{1/2}$, elle vaut $c_0/4$. Pour une droite, les coefficients directeurs des segments successifs sont égaux.\par\smallskip
\textbf{Indice 2.} Si le modèle d'ordre 2 convient, les points $(t\,;\,1/c)$ sont alignés et le coefficient directeur
vaut $k$. Pour $t_{1/2}$, remplacer $c$ par $c_0/2$ dans la relation admise ; pour la division suivante, prendre
$c_0/2$ comme nouvelle concentration initiale.\par\smallskip
\editionEtape{D05 · Une horloge à iode}
\textbf{Indice 1.} Avant $t_\mathrm{c}$, tout le diiode formé a été consommé par les ions thiosulfate : la
stœchiométrie de (2) donne la concentration de diiode produite. La réaction (2) rend les deux ions iodure que (1) a
consommés.\par\smallskip
\textbf{Indice 2.} Si la vitesse est constante, $[\ce{I2}]_\text{produit}=v\,t_\mathrm{c}$, d'où
$k=\dfrac{[\ce{I2}]_\text{produit}}{t_\mathrm{c}\,[\ce{H2O2}]_0\,[\ce{I-}]_0}$ pour chaque expérience. Pour le
mécanisme, additionner les quatre actes et repérer les entités formées puis consommées.\par\smallskip
\editionEtape{D06 · Le générateur de technétium}
\textbf{Indice 1.} Pendant une durée $\mathrm{d}t$, il se forme autant de noyaux de technétium que de noyaux de
molybdène qui se désintègrent, soit $\lambda_1N_1\,\mathrm{d}t$, et il en disparaît $\lambda_2N_2\,\mathrm{d}t$.\par\smallskip
\textbf{Indice 2.} $A_2=\lambda_2N_2$ est maximale lorsque $\frac{\mathrm{d}N_2}{\mathrm{d}t}=0$ : l'équation
différentielle donne alors directement $\lambda_1N_1=\lambda_2N_2$. Pour la limite, écrire $\frac{A_2}{A_1}=\frac{\lambda_2N_2}{\lambda_1N_1}$
et factoriser par $\mathrm{e}^{-\lambda_1t}$.\par\smallskip
\editionEtape{D07 · Le critère d'évolution vu par les vitesses}
\textbf{Indice 1.} L'espèce C est formée par l'acte direct et consommée par l'acte inverse. À l'équilibre, sa
concentration ne varie plus.\par\smallskip
\textbf{Indice 2.} Factoriser $v_1$ dans $v_1-v_{-1}$ et faire apparaître $\frac{k_{-1}}{k_1}=\frac1K$. Pour la composition
à l'équilibre, prendre la racine carrée de l'égalité $Q_{\mathrm{r},\text{éq}}=K$.\par\smallskip
\editionEtape{D08 · Quand une pile est-elle usée ?}
\textbf{Indice 1.} Au départ, $Q_\mathrm{r}=[\ce{Zn^2+}]/[\ce{Cu^2+}]=1$. La tension s'annule si $0{,}0296\log Q_\mathrm{r}=1{,}10$.
Chaque fois que deux électrons circulent, un ion \ce{Cu^2+} disparaît et un ion \ce{Zn^2+} apparaît.\par\smallskip
\textbf{Indice 2.} Lorsque $U=0$, $[\ce{Cu^2+}]=[\ce{Zn^2+}]/K$, avec $[\ce{Zn^2+}]\approx2c$. Le nombre d'ions vaut
$N=n\,N_\mathrm{A}$.\par\smallskip
\end{blocEdition}

\editionPage{indices-defis-3}{Indications progressives (suite)}{Lire une aide, puis reprendre la recherche}{Les indices orientent la démarche ; les solutions restent dans la partie Corrigés.}
\begin{blocEdition}{indices-defis-3}
\editionEtape{D09 · Le pH d'un acide très dilué}
\textbf{Indice 1.} Les charges positives et négatives se compensent : $[\ce{H3O+}]=[\ce{Cl-}]+[\ce{HO-}]$. Remplacer
$[\ce{HO-}]$ par $K_\mathrm{e}/h$ grâce au produit ionique, puis multiplier par $h$.\par\smallskip
\textbf{Indice 2.} Si $c$ est très grand devant $\sqrt{K_\mathrm{e}}$, la racine vaut presque $c$ ; si $c$ est très
petit, elle tend vers $\sqrt{K_\mathrm{e}}=\qty{1,0e-7}{mol.L^{-1}}$, c'est-à-dire vers le pH de l'eau pure.\par\smallskip
\editionEtape{D10 · Pourquoi au moins 1,23 volt ?}
\textbf{Indice 1.} Tant que $U<U_0$, la « pile » formée par les gaz aux électrodes l'emporte sur le générateur : le
courant ne peut pas circuler dans le sens de l'électrolyse. L'énergie électrique reçue vaut $W=Q\,U$, et chaque mole de
dihydrogène échange deux moles d'électrons.\par\smallskip
\textbf{Indice 2.} Pour la dernière question, comparer 237~kJ (reçus du générateur) et 242~kJ (restitués par la
combustion) : le complément ne peut venir que d'un transfert thermique. La tension « thermoneutre » vérifie
$2F\,U=\qty{242}{kJ}$.\par\smallskip
\editionEtape{D11 · Combien d'isomères ?}
\textbf{Indice 1.} Classer les squelettes selon la longueur de leur chaîne la plus longue : pour \ce{C5H12}, elle vaut
5, 4 ou 3. Un groupe méthyle placé au bout d'une chaîne ne crée pas de ramification : il l'allonge.\par\smallskip
\textbf{Indice 2.} Nommer chaque squelette obtenu : deux dessins qui reçoivent le même nom sont le même isomère. Pour
\ce{C6H14}, essayer de placer un groupe éthyle sur une chaîne de quatre carbones et nommer le résultat.\par\smallskip
\editionEtape{D12 · Jusqu'où pousser un rendement ?}
\textbf{Indice 1.} À l'équilibre, $[\text{ester}][\text{eau}]=4\,[\text{acide}][\text{alcool}]$ : avec 1~mol d'acide et
$r$~mol d'alcool, les quantités finales valent $x$, $x$, $1-x$ et $r-x$. Développer l'égalité $x^2=4(1-x)(r-x)$.\par\smallskip
\textbf{Indice 2.} Garder la racine inférieure à 1 du trinôme. Pour la dernière question : retirer l'eau maintient
$Q_\mathrm{r}<K$ (le système n'atteint jamais le palier), alors qu'un catalyseur ne change pas $K$.\par\smallskip
\end{blocEdition}

\editionPage{indices-defis-4}{Indications progressives (suite)}{Lire une aide, puis reprendre la recherche}{Les indices orientent la démarche ; les solutions restent dans la partie Corrigés.}
\begin{blocEdition}{indices-defis-4}
\editionEtape{D13 · Le caillou coincé dans un pneu}
\textbf{Indice 1.} Dériver $\sin(\omega t)$ donne $\omega\cos(\omega t)$ et dériver $\cos(\omega t)$ donne
$-\omega\sin(\omega t)$. Le caillou touche le sol quand $\omega t$ est un multiple de $2\pi$.\par\smallskip
\textbf{Indice 2.} Comparer les coordonnées de $\vec{a}$ à celles du vecteur $\vect{MC}=(x_\mathrm{C}-x\,;\,y_\mathrm{C}-y)$ :
elles sont proportionnelles. Dans le référentiel du vélo, le centre de la roue est immobile.\par\smallskip
\editionEtape{D14 · De quelle hauteur lancer le wagon ?}
\textbf{Indice 1.} Entre le départ et le sommet, le wagon descend de $h-2R$ : $\frac12\,m\,v^2=m\,g\,(h-2R)$. Au sommet,
$\vec{P}$ et $\vec{N}$ sont tous deux dirigés vers le bas, c'est-à-dire selon $\un$.\par\smallskip
\textbf{Indice 2.} Au point bas, $\vec{N}$ est dirigé vers le haut (vers le centre) et $\vec{P}$ vers le bas : selon
$\un$, $N-m\,g=m\,v^2/R$ avec $v^2=2\,g\,h$.\par\smallskip
\editionEtape{D15 · Lancer à 45° ?}
\textbf{Indice 1.} Depuis le sol, $z=0$ pour $t=\dfrac{2\,v_0\sin\alpha}{g}$, puis utiliser $2\sin\alpha\cos\alpha=\sin(2\alpha)$.
Depuis la hauteur $h$, résoudre $-\frac12\,g\,t^2+v_z\,t+h=0$ et garder la racine positive.\par\smallskip
\textbf{Indice 2.} Pendant la chute supplémentaire de $h$, le projectile avance horizontalement : cette phase profite aux
lancers rasants, dont la vitesse horizontale est plus grande.\par\smallskip
\editionEtape{D16 · Peser le trou noir de la Galaxie}
\textbf{Indice 1.} Écrire $\dfrac{T^2}{a^3}=\dfrac{4\pi^2}{GM}$ pour l'étoile et $\dfrac{T_\oplus^2}{a_\oplus^3}=\dfrac{4\pi^2}{GM_\odot}$ pour la Terre,
puis diviser membre à membre : $G$ et $4\pi^2$ disparaissent.\par\smallskip
\textbf{Indice 2.} Distance minimale : $a(1-e)$. Loi des aires aux extrémités : $r_\text{min}\,v_\text{max}=r_\text{max}\,v_\text{min}$,
comme pour la comète de Halley (exercice~6).\par\smallskip
\editionEtape{D17 · La clepsydre}
\textbf{Indice 1.} Le volume perdu par le vase pendant $\mathrm{d}t$ vaut $\pi\,r^2\,|\mathrm{d}h|$ ; le volume sorti vaut
$s\,v\,\mathrm{d}t$. Pour le cylindre, dériver $\sqrt{h}$ donne $\dfrac{1}{2\sqrt{h}}\,\dfrac{\mathrm{d}h}{\mathrm{d}t}$.\par\smallskip
\textbf{Indice 2.} Imposer $u$ constant dans $\pi\,r^2\,u=s\,\sqrt{2\,g\,h}$ : $r^2$ est proportionnel à $\sqrt{h}$. Au sommet,
$r=R$ pour $h=H$. Pour la question~5, exprimer $s$ à l'aide de $u$ dans la durée de vidange du cylindre.\par\smallskip
\end{blocEdition}

\editionPage{indices-defis-5}{Indications progressives (suite)}{Lire une aide, puis reprendre la recherche}{Les indices orientent la démarche ; les solutions restent dans la partie Corrigés.}
\begin{blocEdition}{indices-defis-5}
\editionEtape{D18 · Faire le vide}
\textbf{Indice 1.} Pendant l'aspiration, l'air du récipient occupe le volume $V+v$ à température constante :
$p_1\,(V+v)=p_0\,V$. Une suite qui vérifie $p_{k+1}=q\,p_k$ est géométrique de raison $q$.\par\smallskip
\textbf{Indice 2.} Résoudre $p_0\,q^k<\qty{1}{hPa}$ avec le logarithme népérien, sans oublier que $\ln q<0$. Avec un volume
mort, la pression limite est la valeur de $p$ qui ne change plus d'un coup au suivant.\par\smallskip
\editionEtape{D19 · L'expérience de Joule}
\textbf{Indice 1.} Aucun corps plus chaud ne touche l'eau : le système ne reçoit que du travail. Pour une descente des deux
masses, ce travail vaut $m\,g\,h-\frac12\,m\,v^2$, l'énergie mécanique qu'elles perdent.\par\smallskip
\textbf{Indice 2.} La « valeur en eau » du récipient est la masse d'eau de même capacité thermique : $C/c_\mathrm{eau}$.
Exprimer l'énergie reçue en calories avec la masse d'eau équivalente totale, puis diviser le travail par ce nombre.\par\smallskip
\editionEtape{D20 · Café : le lait avant ou après ?}
\textbf{Indice 1.} Pendant le mélange, l'énergie reçue par le lait est égale à l'énergie cédée par le café. Pendant l'attente,
la boisson suit $T(t)=T_\mathrm{e}+(T_0-T_\mathrm{e})\,\mathrm{e}^{-t/\tau}$, où $T_0$ est sa température au début de l'attente.\par\smallskip
\textbf{Indice 2.} Pour la formule générale, développer $T_\mathrm{A}$ et $T_\mathrm{B}$ avec $k=\mathrm{e}^{-t/\tau}$ : les termes qui
contiennent $T_\mathrm{c}$ se compensent.\par\smallskip
\editionEtape{D21 · Pourquoi entend-on les basses du voisin ?}
\textbf{Indice 1.} $\log(m_\mathrm{s}\times2f)=\log(m_\mathrm{s}\,f)+\log2$. Le niveau transmis de l'autre côté de la paroi vaut
$L-R$.\par\smallskip
\textbf{Indice 2.} Pour la porte, calculer la fraction transmise par chaque partie, en faire la moyenne pondérée par les
surfaces, puis revenir à un indice avec $R=-10\log\tau$.\par\smallskip
\editionEtape{D22 · Des franges dans la tache de diffraction}
\textbf{Indice 1.} La tache centrale de diffraction s'étend de $-\lambda D/a$ à $\lambda D/a$. Au centre d'une frange brillante, le
terme d'interférences vaut 1.\par\smallskip
\textbf{Indice 2.} Compter les entiers $k$ tels que $|k\,i|<\dfrac{\lambda D}{a}$, en excluant la frange qui tombe exactement sur un
zéro de diffraction.\par\smallskip
\end{blocEdition}

\editionPage{indices-defis-6}{Indications progressives (suite)}{Lire une aide, puis reprendre la recherche}{Les indices orientent la démarche ; les solutions restent dans la partie Corrigés.}
\begin{blocEdition}{indices-defis-6}
\editionEtape{D23 · Le passage d'une voiture}
\textbf{Indice 1.} Quand $x\to-\infty$, $\dfrac{-x}{\sqrt{x^2+b^2}}\to1$. Le son émis à la date $t$ parcourt la distance $d$ à la célérité
$v$ avant d'être reçu.\par\smallskip
\textbf{Indice 2.} Éliminer $f_\mathrm{E}$ entre les expressions des deux fréquences extrêmes. La fréquence change surtout quand
$|x|$ est de l'ordre de $b$.\par\smallskip
\editionEtape{D24 · Où placer l'œil ? Le cercle oculaire}
\textbf{Indice 1.} Pour l'oculaire, l'objet est l'objectif, en $\overline{OA}=-(f'_1+f'_2)$ ; réduire au même dénominateur après avoir
appliqué la relation de conjugaison.\par\smallskip
\textbf{Indice 2.} Le grandissement donne le rapport des diamètres de l'image et de l'objectif. Toute la lumière entre dans
l'œil si le cercle oculaire est plus petit que la pupille.\par\smallskip
\editionEtape{D25 · Mesurer la constante de Planck avec des LED}
\textbf{Indice 1.} L'énergie électrique reçue par un électron se retrouve dans le photon émis. Une relation de la forme
$U_\mathrm{s}=a\times\dfrac{1}{\lambda}$ représente une droite passant par l'origine.\par\smallskip
\textbf{Indice 2.} $h=\dfrac{a\,e}{c}$. Pour comparer chaque mesure au modèle, calculer $\dfrac{1239{,}8}{\lambda}$ avec $\lambda$ en nanomètres.\par\smallskip
\editionEtape{D26 · Où passe la moitié de l'énergie ?}
\textbf{Indice 1.} Sous une tension constante $E$, l'énergie fournie par le générateur vaut $E\,q$. Comparer à l'énergie stockée
$\frac12\,C\,E^2$.\par\smallskip
\textbf{Indice 2.} Une primitive de $\mathrm{e}^{-2t/\tau}$ est $-\dfrac{\tau}{2}\,\mathrm{e}^{-2t/\tau}$. Dans une charge par étapes, la tension aux bornes
de la résistance part, à chaque étape, de $E/N$.\par\smallskip
\editionEtape{D27 · Mesurer plus pour mesurer mieux ?}
\textbf{Indice 1.} Écrire $\bar m=\dfrac{1}{N}\left(m_1+\dots+m_N\right)$ : la somme a pour incertitude-type $\sqrt N\,s$. Un retard qui
s'ajoute à chaque mesure se retrouve tel quel dans la moyenne.\par\smallskip
\textbf{Indice 2.} L'écart normalisé $z=\dfrac{\delta}{s/\sqrt N}$ croît comme $\sqrt N$. Pour la dernière question, tracer $t_\mathrm{mes}$ en
fonction de $\sqrt h$ et chercher une droite.\par\smallskip
\end{blocEdition}
